\scanpage{109}
\mspage{3.2}
\uline{Proof:} 1) Assume that $S$ is finitely generated as $R$-module. Then $\mathrm{w.\ hd}_R S<\infty$ implies $\mathrm{hd}_R S<\infty$. Let $0\longrightarrow X_t\longrightarrow\cdots\longrightarrow X_1\longrightarrow X_0\longrightarrow S\longrightarrow0$ be an acyclic resolution of $S$ by $R$-free modules of finite rank, and let $\underline m$ be the maximal ideal of $R$. Then by the additivity of the multiplicity, we have
\[
\begin{aligned}
\frac{e_S(\underline m)}{e_R(\underline m)}
&=\sum_i(-1)^i e_{X_i}(\underline m)/e_R(\underline m)=\sum_i(-1)^i[X_i:R]\\
&=\sum_i(-1)^i[X_i\otimes_R\overline R:R]=\sum_i(-1)^i[\operatorname{Tor}_i^R(S,\overline R):\overline R]\\
&=[\overline S:\overline R]\cdot[\sum_i(-1)^iL_S(\operatorname{Tor}_i^R(S,\overline R))].
\end{aligned}
\]
Therefore
\[
e(S/R)=\frac{e_S(\underline m S)}{e_R(\underline m)}=\frac1{[\overline S:\overline R]}\cdot\frac{e_S(\underline m)}{e_R(\underline m)}=\sum_i(-1)^iL_S(\operatorname{Tor}_i^R(S,\overline R)).
\]

2) General case: Since $\overline S$ is finite dimensional over $\overline R$, $\widehat S$ is finitely generated as $\widehat R$-module. Furthermore, by lemma 3.1, we have $\mathrm{hd}_{\widehat R}\widehat S<\infty$. Therefore we are brought to the above situation and hence
\[
e(S/R)=e(\widehat S/\widehat R)=\sum_i(-1)^iL_{\widehat S}(\operatorname{Tor}_i^{\widehat R}(\widehat S,\overline R))=\sum_i(-1)^iL_S(\operatorname{Tor}_i^R(S,\overline R)).
\]

Corollary 3.3: Let $S$ be a local $R$-algebra over a local ring $R$. If $\mathrm{w.\ hd}_S S^f<\infty$, then $r(S/R)$ is an integer, equal to
\[
\sum_i(-1)^iL_{S^f}(\operatorname{Tor}_i^S(S^f\underline m\,\overline S))=\sum_i(-1)^iL_{S^f}(S^f\otimes_{S^e}\operatorname{Tor}_i^S(S^e,\overline S)).
\]

Lemma 3.4: Let $R\subset S$ be local domains such that $S$ is finitely generated as $R$-module. If $R$ is integrally closed and $\overline S$ is a simple extension of $\overline R$, then $e(S/R)$ is an integer.

\uline{Proof:} Let $\alpha$ be an element in $S$ such that $\overline S=\overline R(\overline\alpha)$ where $\overline\alpha$ is the image of $\alpha$ under the canonical map $S\longrightarrow\overline S$. The subring $R[\alpha]$ is $R$-free since $R$ is integrally closed, and is a local $R$-algebra since $S$ is a finite local $R$-algebra. Therefore $e(R[\alpha]/R)=[R[\alpha]:R]/[\overline R[\overline\alpha]:\overline R]=[R[\alpha]:R]/[\overline S:\overline R]$ is an integer by proposition 3.2. Consequently $e(S/R)=[S:R]/[\overline S:\overline R]=[S:R[\alpha]]\cdot e(R[\alpha]/R)$ is an integer.
\mspage{108}
